\chapter{Conclusions and future work}

\section{Conclusions}

We developed the necessary computational tools for the Finite Element modeling, simulation and analysis of EM wave propagation in the context of Photonic Crystals.  This was done as an ongoing open source software platform based on Python and structured around Object Oriented Paradigm. 
In order to build this platform, the necessary concepts and mathematical tools associated to the modeling of electromagnetic waves in periodic media were appropriated, specifically Maxwell's equations, and periodicity formulations taken from solid state physics. 

Moreover, the finite element method was explored as an alternative to other methods commonly used in the analysis of PCs such as the Plane Wave Expansion (PWE) \cite{StevenJohnson2001} method\add[SEC]{,} and Finite Differences Time Domain (FDTD)\cite{Oskooi2009}. FEM was chosen mainly because it has proven to be a great general purpose method for investigating PCs \cite{Andonegui2013} being its flexibility to model interfaces and defects in both frequency and spatial domains one of its greatest advantages. FEM is also well known for its stability and robustness and has proven its versatility in many fields of engineering \cite{Bathe1996, Ram2002, Zienkiewicz2005, Logg2012}.

The platform was designed using the Object Oriented Paradigm having in mind the objective of \change[SEC]{having}{implementing} a well structured and flexible architecture, capable of being upgraded to future capabilities. \remove{Even though the architecture is in a developing stage and a proper manual for its use its pending,} The modules, classes and methods that conform it are functional for production of results, and are documented for user reference. The reader is encouraged to visit the \href{https://github.com/bebopsan/peyeQM/tree/Depuration}{\textbf{online repository}} at \url{https://github.com/bebopsan/peyeQM/tree/Depuration} and explore its examples. The documentation page is also available online at: \url{http://bebopsan.github.io/peyeQM}.

We proved that PeyeQM is capable of accurately simulating  2D electromagnetism stationary, and harmonic vector fields as well as fields inside in-homogeneous domains and infinite crystals. The need of CAD tools for the construction of finite lattices of dielectric rods with arbitrary size and placement was a factor that was underestimated in the beginning, but the outcome ended up becoming a great addition that gives flexibility to the modeling of complex domains.

\change[SEC]{In the} At present time, we are working on fixing bugs, cleaning the code, and making a detailed documentation of the software in order to achieve a product that can be shared and published as an article and or a downloadable application. This is a work that has great potential for academic and educational use by students of engineering and physics interested in diving into computational physics in the context of open source collaborative software.  
 

\section{Future work}
The following is a list of items that are worth exploring in future work on PeyeQM platform:
\begin{itemize}
%
\item Definition of vectorial triangular elements for improving meshing flexibility and linear quadrilateral elements for computation speed. Triangular elements are simpler than quadrilaterals and are very often used by meshing algorithms in order to fill regions where quadrilateral elements get distorted. On the other hand, a comparison between convergence and speed between linear and quadratic elements would improve the understanding of the problem and ease calculations.
%
\item Consider adding support for infinite boundary conditions for the simulation of unbounded domains \cite{antoine2009review, appelo2003}. One popular candidate very often used in electromagnetism is the use of regions with Perfectly Matched Layers (PML)\cite{Jin2010}, where the domain is surrounded with a material that has artificial energy dissipation and the properties are fitted to minimize reflection between the interfaces. Other possibilities are:
\begin{itemize}
\item to use a Robin boundary condition when intrinsic impedance is known, or to map the unbounded region to a known geometry (e.g, with a conformal mapping), both approaches used in the open source code FEMM (\href{http://www.femm.info/wiki/HomePage}{\textbf{Finite Element Method Magnetics}}) \cite{meeker2010_FEMM},
\item implement hybrid FEM and Boundary Elements(BEM) as presented by Nicol\'as Guar\'in in \cite{Guarin2012},
\item or define infinite Elements, that are special finite elements with an assumed decaying behavior (\cite{Zienkiewicz2005}).
\end{itemize} 
%
\item Evaluate convergence of the method for the problems described in chapter \ref{ch:Results} using finer spatial and spectral ($\vec{k}$ domain) meshes.
%
\item Perform simulations of time dependent problems involving defects in finite crystals, and compare efficiency of designs involving PC waveguides or cavities with publications.
%
\item Integrate the CAD and solver into an optimization algorithm that maximizes transmission or confinement by inducing variations in position and shape of defects in the lattice.
%
\item Consider parallelization schemes to be implemented in order to make bigger simulations plausible. Algorithms that split the computation into threads and take advantage of the parallel architecture of modern computers are the key to obtaining results of enormous simulation faster than before. Making research on things like weather and seismic simulation, molecular dynamics or drug compounds testing practical and realizable. With a parallel implementation of PeyeQM bigger domains and finer meshes would be able to be modelled in a robust workstation. It is to notice that nowadays parallel computation capacity is available for rent in places such as Amazon's \href{http://aws.amazon.com/ec2/}{ec2} service.
\item Possible integration of (now external) processes such as meshing and visualization, by means of scripting that uses gmsh or Paraview Python libraries. This would make the process of solving more autonomous and free the user from external tasks.
\end{itemize}

\chapter{Software and document licenses}

\section{Thesis document license}
\begin{figure}[h!]
\includegraphics[scale=2]{./img/CC-BY-SA_icon.pdf}
\end{figure}
All the contents of this work are shared by a Creative Commons Atribución-CompartirIgual 2.5 Colombia. Contents that do not belong to the author can not be used  under the terms of this license.
The user is encouraged to read about the reach and conditions of the license in the following web pages: 
\begin{itemize}
\item \url{http://creativecommons.org/licenses/by-sa/2.5/co/}
\item \url{http://creativecommons.org/licenses/by-sa/2.5/co/legalcode}
\end{itemize} 
\section{PeyeQM license}

The software developed under this project is distributed under the \textbf{standard MIT Licence},  please use it and twist it in whatever fashion you wish.

Copyright (C) 2013 Santiago Echeverri Chacón, Nicolás Guarín Zapata, Edward Yesid Villegas Pulgarín.

Permission is hereby granted, free of charge, to any person obtaining a copy of this software and associated documentation files (the "Software"), to deal in the Software without restriction, including without limitation the rights to use, copy, modify, merge, publish, distribute, sublicense, and/or sell copies of the Software, and to permit persons to whom the Software is furnished to do so, subject to the following conditions:

The above copyright notice and this permission notice shall be included in all copies or substantial portions of the Software.

THE SOFTWARE IS PROVIDED "AS IS", WITHOUT WARRANTY OF ANY KIND, EXPRESS OR IMPLIED, INCLUDING BUT NOT LIMITED TO THE WARRANTIES OF MERCHANTABILITY, FITNESS FOR A PARTICULAR PURPOSE AND NONINFRINGEMENT. IN NO EVENT SHALL THE AUTHORS OR COPYRIGHT HOLDERS BE LIABLE FOR ANY CLAIM, DAMAGES OR OTHER LIABILITY, WHETHER IN AN ACTION OF CONTRACT, TORT OR OTHERWISE, ARISING FROM, OUT OF OR IN CONNECTION WITH THE SOFTWARE OR THE USE OR OTHER DEALINGS IN THE SOFTWARE.

